\documentclass[12pt]{article}
\usepackage{amsmath,amssymb,amsthm}
\usepackage[margin=1in]{geometry}

\title{Problem 267: Billionaire}
\author{Project Euler Solution}
\date{}

\begin{document}
\maketitle

\section{Problem Statement}
You start with \$1. You toss a fair coin 1000 times. Before each toss, you choose a fixed fraction $f \in (0,1)$ of your current wealth to bet. If heads, you gain $2f$ times your wealth (multiply by $1+2f$); if tails, you lose $f$ times your wealth (multiply by $1-f$). Find the value of $f$ that maximizes the probability of having at least \$1,000,000,000 after 1000 tosses. Give this probability to 12 decimal places.

\section{Mathematical Analysis}

\subsection{Wealth Model}
After $n = 1000$ tosses with $k$ heads and $n-k$ tails:
\[ W(k) = (1+2f)^k \cdot (1-f)^{n-k} \]

We need $W(k) \geq 10^9$, i.e.:
\[ k \cdot \ln(1+2f) + (n-k) \cdot \ln(1-f) \geq 9\ln 10 \]

\subsection{Optimal Fraction for Given $k$}
For a fixed number of heads $k$, the wealth $W(k)$ is maximized over $f$ at:
\[ \frac{\partial}{\partial f}\left[k\ln(1+2f) + (n-k)\ln(1-f)\right] = \frac{2k}{1+2f} - \frac{n-k}{1-f} = 0 \]
Solving: $f^* = \frac{3k - n}{2n}$, valid when $k > n/3$.

\subsection{Minimum Heads Required}
Substituting $f^*$, the maximum wealth for $k$ heads is:
\[ W_{\max}(k) = \left(\frac{3k}{n}\right)^k \cdot \left(\frac{3(n-k)}{2n}\right)^{n-k} \]

The minimum $k$ (call it $k_{\min}$) satisfying $W_{\max}(k) \geq 10^9$ is found numerically. For $n = 1000$, $k_{\min} = 432$.

\subsection{Probability Calculation}
The probability of at least $k_{\min}$ heads out of $n$ fair coin tosses:
\[ P = \sum_{k=k_{\min}}^{n} \binom{n}{k} \cdot 2^{-n} = \frac{1}{2^{1000}} \sum_{k=432}^{1000} \binom{1000}{k} \]

This is computed exactly using integer arithmetic for $\binom{1000}{k}$ and arbitrary precision division.

\subsection{Proof of Optimality}
For any $f$, the minimum number of heads needed is:
\[ k_{\min}(f) = \left\lceil \frac{9\ln 10 - n\ln(1-f)}{\ln(1+2f) - \ln(1-f)} \right\rceil \]

The continuous function $g(f)$ being minimized achieves its minimum at the $f^*$ corresponding to the boundary $k = k_{\min}$. Since the probability $P$ depends only on $k_{\min}$ (through the binomial CDF), the optimal $f$ is any value achieving the minimum $k_{\min}$.

\section{Complexity}
\begin{itemize}
    \item Finding $k_{\min}$: $O(n)$ evaluations of logarithmic expressions.
    \item Computing probability: $O(n)$ binomial coefficient computations.
    \item Total: $O(n)$ where $n = 1000$.
\end{itemize}

\section{Answer}

\[
\boxed{0.999992836187}
\]

\end{document}
